%%%%%%%%%%%%%%%%%%%%%%%%%%%%%%%%%%%%%%%%%%%%%%%%%%%%%%%%%%%%%%%%%%%%%%
% Lớp 9 | Chương V. Đường tròn | Độ dài của cung tròn. Diện tích hình quạt tròn và hình vành khuyên
% Dạng 9H5?15-3: Diện tích hình vành khuyên, bài toán thực tế
%%%%%%%%%%%%%%%%%%%%%%%%%%%%%%%%%%%%%%%%%%%%%%%%%%%%%%%%%%%%%%%%%%%%%%


% nguồn: dan_tu_AI (nhập 14/11/2026)
\begin{ex}%[9H5N15-3]%[ma:473df6a6]
	Diện tích $S$ của hình vành khuyên giới hạn bởi hai đường tròn đồng tâm $(O;R)$ và $(O;r)$ với $R>r$ là
	\choice
	{$S=\pi(R-r)^2$}
	{\True $S=\pi\left(R^2-r^2\right)$}
	{$S=\pi(R-r)$}
	{$S=2\pi(R-r)$}
	\loigiai{
		$S=\pi R^2-\pi r^2=\pi\left(R^2-r^2\right)$.
	}
\end{ex}

% nguồn: dan_tu_AI (nhập 14/11/2026)
\begin{ex}%[9H5H15-3]%[ma:23f68989]
	Diện tích của hình vành khuyên giới hạn bởi hai đường tròn đồng tâm có bán kính $5\text{ cm}$ và $3\text{ cm}$ là
	\choice
	{$4\pi\text{ cm}^2$}
	{$2\pi\text{ cm}^2$}
	{$34\pi\text{ cm}^2$}
	{\True $16\pi\text{ cm}^2$}
	\loigiai{
		$S=\pi\left(5^2-3^2\right)=\pi(25-9)=16\pi\text{ (cm}^2)$.
	}
\end{ex}

% nguồn: dan_tu_AI (nhập 14/11/2026)
\begin{bt}%[9H5V15-3]%[ma:22909110]
	Một hồ nước hình tròn có bán kính $4\text{ m}$. Người ta làm một lối đi lát gạch xung quanh hồ, rộng $1{,}5\text{ m}$ (phần lối đi là hình vành khuyên). Tính diện tích lối đi (lấy $\pi\approx 3{,}14$, đơn vị m$^2$, làm tròn đến hàng phần mười).
	\shortans{44,7}
	\loigiai{
		Bán kính đường tròn ngoài: $R=4+1{,}5=5{,}5\text{ (m)}$.\\
		\[ S=\pi\left(R^2-r^2\right)=\pi\left(5{,}5^2-4^2\right)=14{,}25\pi\approx 44{,}745\approx 44{,}7\text{ (m}^2). \]
	}
\end{bt}
